
\documentclass[11pt,a4paper]{article}
\usepackage[T1]{fontenc}
\usepackage[utf8]{inputenc}
\usepackage[ngerman]{babel}
\usepackage{lmodern}
\usepackage{microtype}
\usepackage{amsmath,amssymb,amsthm,mathtools}
\usepackage{booktabs,tabularx,array}
\usepackage{geometry}
\usepackage{xcolor}
\usepackage{enumitem}
\usepackage{hyperref}
\usepackage{tcolorbox}
\geometry{left=24mm,right=24mm,top=23mm,bottom=25mm}
\definecolor{deepblue}{HTML}{1F4E79}
\definecolor{deepred}{HTML}{8B1E1E}
\definecolor{softgray}{HTML}{F3F4F6}
\hypersetup{
  colorlinks=true,
  linkcolor=deepblue,
  citecolor=deepblue,
  urlcolor=deepblue,
  pdftitle={EABC Forschungsbericht v23 - Kapitel 7: Yang-Mills Bridge Problem}
}
\setlength{\parindent}{0pt}
\setlength{\parskip}{0.55em}
\newtheorem{satz}{Satz}
\newtheorem{lemma}{Lemma}
\newtheorem{definition}{Definition}
\newtheorem{annahme}{Annahme}
\newtheorem{vermutung}{Brückenvermutung}
\newtcolorbox{statusbox}[1]{colback=softgray,colframe=deepblue,title=\textbf{#1},boxrule=.7pt,arc=1mm}
\newtcolorbox{warningbox}{colback=white,colframe=deepred,title=\textbf{Methodische Grenze},boxrule=.8pt,arc=1mm}

\begin{document}

\begin{center}
{\Large\bfseries EABC Forschungsbericht v23}\\[2mm]
{\large Kapitel 7: Vom algebraischen Skelett zur analytischen Yang--Mills-Brücke}\\[1mm]
{\small Revidierte, epistemisch getrennte Fassung}
\end{center}

\begin{warningbox}
Die in diesem Kapitel verwendeten gruppentheoretischen, kohomologischen,
VOA-theoretischen und oktonionischen Strukturen liefern Bedingungen für die
\emph{interne algebraische Konsistenz} des EABC-Modells. Insbesondere implizieren weder
\[
[\alpha]=0
\]
für eine relevante $H^3$-Anomalieklasse noch
\[
P_{\rm phys}[x,y,z]_{\mathbb O}=0
\]
für den projizierten oktonionischen Assoziator die Existenz einer positiven
Massenlücke einer vierdimensionalen quantisierten Yang--Mills-Theorie.

Diese Bedingungen kontrollieren Lift-, Assoziativitäts- und Sektorobstruktionen.
Sie enthalten für sich genommen keine quantitative Information über die untere
Spektralschranke eines vierdimensionalen Yang--Mills-Hamiltonoperators.

Daher gilt im gesamten Kapitel strikt:
\[
\boxed{
\text{algebraische Anomaliefreiheit}
\;\not\Rightarrow\;
\text{analytische Massenlücke}.
}
\]
Ebenso gilt:
\[
\boxed{
\text{Nullmoden- oder Leakage-Kontrolle}
\;\not\Rightarrow\;
\Delta_{\rm YM}>0.
}
\]
\end{warningbox}

\section*{7.1 Logische Verortung der algebraischen Strukturen}

Die bisher untersuchten Strukturen
\[
V_4,\qquad Q_8,\qquad H^2,\qquad H^3,\qquad
\text{VOA/Leech},\qquad \mathbb O
\]
stehen logisch \emph{vor} der analytischen Yang--Mills-Kette. Sie können einen
diskreten Zustandsraum organisieren, projektive Liftklassen kodieren,
kategorielle Kohärenzbedingungen formulieren und bestimmte interne
Obstruktionen ausschließen.

Ihre mathematische Rolle ist damit kinematisch beziehungsweise algebraisch.
Sie liefern nicht automatisch Lorentz-Kovarianz, eine vierdimensionale
Kontinuumstheorie oder eine positive untere Spektralschranke.

\begin{statusbox}{Epistemische Dreiteilung}
\[
\boxed{
\begin{array}{c}
\text{algebraische / kategorielle Kohärenz}\\
\Downarrow\\
\text{regulierte spektrale Modellstruktur}\\
\Downarrow\\
\text{offene analytische Kontinuumsbrücke}
\end{array}}
\]
Kein Pfeil in dieser Kette wird allein durch Symmetrie oder Kohomologie ersetzt.
\end{statusbox}

\section*{7.2 Erste Hürde: eichinvariante vierdimensionale Quantenstruktur}

Ein endlicher oder diskreter Gitter-Hamiltonoperator kann ein wohldefiniertes
Spektralmodell liefern, ist aber noch keine quantisierte Yang--Mills-Theorie auf
$\mathbb R^4$.

Für eine Clay-relevante Brücke muss eine mathematisch kontrollierte
vierdimensionale Theorie konstruiert werden. In einer euklidischen Formulierung
bedeutet dies schematisch eine Familie von eichinvarianten Korrelationsfunktionen
oder ein geeignetes Maß auf regulierten Eichfeldkonfigurationen, aus dem eine
physikalische Theorie rekonstruiert werden kann.

Die logische Zielkette lautet:
\[
\boxed{
\text{regulierte Eichfelder}
\longrightarrow
\text{euklidische QFT}
\longrightarrow
\text{physikalischer Hilbertraum}.
}
\]

Eichinvarianz, Lokalität, Nichttrivialität und die geeigneten
Rekonstruktionsaxiome müssen im Grenzübergang erhalten bleiben.

\section*{7.3 Zweite Hürde: kontrollierter UV-/IR-Grenzübergang}

Sei
\[
H_{a,L}
\]
eine regulierte Familie, wobei $a>0$ den Gitterabstand und $L<\infty$ die
Volumenskala bezeichnet.

Die Kontinuumstheorie erfordert
\[
\boxed{
a\to0,\qquad L\to\infty.
}
\]

Der erste Grenzübergang entfernt den UV-Regulator, der zweite den
endlichen Volumenregulator. Es genügt nicht, dass jedes einzelne regulierte
Modell mathematisch wohldefiniert ist.

Für geeignete Observablen $\mathcal O_1,\ldots,\mathcal O_n$ muss vielmehr eine
kontrollierte Konvergenz der Form
\[
\lim_{\substack{a\to0\\L\to\infty}}
\left\langle
\mathcal O_1\cdots\mathcal O_n
\right\rangle_{a,L}
=
S_n
\]
zu einer konsistenten Familie von Kontinuumsfunktionen $S_n$ nachgewiesen
werden.

\section*{7.4 Dritte Hürde: uniforme Spektrallücke}

Für jedes regulierte Modell definiere
\[
\Delta(a,L)
=
\inf\Bigl(
\operatorname{spec}H_{a,L}\setminus\{E_0(a,L)\}
\Bigr)
-
E_0(a,L).
\]

Die Aussage
\[
\Delta(a,L)>0
\]
für jedes feste $a,L$ ist nicht ausreichend. Beispielsweise könnten
\[
\Delta(a,L)\sim a
\qquad\text{oder}\qquad
\Delta(a,L)\sim\frac{1}{\log L}
\]
gelten, sodass die Lücke im physikalischen Grenzübergang verschwindet.

Das analytische Ziel ist eine regulatorunabhängige Koerzivitätsabschätzung:
\[
\boxed{
\exists\,\Delta_*>0:\qquad
\mathfrak q_{a,L}[u]
-
E_0(a,L)\|u\|^2
\ge
\Delta_*\|u\|^2
}
\]
für alle physikalischen Zustände
\[
u\perp\psi_{0,a,L}
\]
und für alle hinreichend kleinen $a$ sowie hinreichend großen $L$.

Insbesondere muss gelten:
\[
\boxed{
\liminf_{\substack{a\to0\\L\to\infty}}
\Delta(a,L)>0.
}
\]

Dies ist die zentrale funktionalanalytische Beweislast. Eine endliche
Gruppen-, Gitter- oder VOA-Symmetrie liefert diese Uniformität nicht von selbst.

\section*{7.5 Vierte Hürde: Osterwalder--Schrader-Rekonstruktion}

Selbst konvergierende euklidische Korrelationsfunktionen ergeben nicht
automatisch eine physikalische relativistische Quantentheorie.

Eine zentrale Positivitätsbedingung ist die Reflexionspositivität. Für geeignete
Observablen $F$, deren euklidischer Zeitträger in der positiven Halbzeit liegt,
muss schematisch
\[
\boxed{
\langle \Theta F,F\rangle_{\rm E}\ge0
}
\]
gelten, wobei $\Theta$ die euklidische Zeitreflexion bezeichnet.

Diese Positivität ist wesentlich für die Konstruktion eines
positiven Hilbertraum-Skalarprodukts in der Osterwalder--Schrader-Rekonstruktion.
Ohne sie ist die gewünschte Rekonstruktion einer unitären physikalischen
Theorie mit positivem Zustandsraum nicht gewährleistet.

Reflexionspositivität ist daher kein kosmetisches Zusatzaxiom, sondern ein
eigenständiges analytisches Gate.


\section*{7.6 Methodische Statushierarchie}

Die relevanten Aussageebenen werden bewusst getrennt:

\begin{itemize}[leftmargin=1.6em]
\item \textbf{\([\alpha]=0\):}
Trivialität einer relevanten kategoriellen Anomalieklasse.
\textbf{Status:} algebraische/kategorielle Bedingung.

\item \textbf{\(P_{\rm phys}\operatorname{Assoc}_{\mathbb O}=0\):}
Neutralität des projizierten oktonionischen Assoziators.
\textbf{Status:} modellinterne Bedingung; separat zu beweisen.

\item \textbf{\(\Delta(a,L)>0\):}
Positive Spektrallücke eines regulierten endlichen Modells.
\textbf{Status:} finite Spektraltheorie.

\item \textbf{\(\liminf \Delta(a,L)>0\):}
Regulatorunabhängige positive Kontinuumslücke.
\textbf{Status:} offen.

\item \textbf{OS-Rekonstruktion:}
Rekonstruktion eines positiven physikalischen Hilbertraums aus der
euklidischen Theorie.
\textbf{Status:} offen.

\item \textbf{4D Yang--Mills + Mass Gap:}
Konstruktive vierdimensionale QFT mit positiver physikalischer Massenlücke.
\textbf{Status:} nicht bewiesen / nicht behauptet.
\end{itemize}

In Kurzform:
\[
\boxed{
[\alpha]=0
\;\neq\;
\Delta(a,L)>0
\;\neq\;
\Delta_{\rm YM}>0.
}
\]

\section*{7.7 Zentrales EABC--Yang--Mills Bridge Theorem Candidate}

Der wissenschaftlich relevante offene Pfeil lautet nicht
\[
\text{Symmetrie}\Longrightarrow\text{Mass Gap},
\]
sondern verlangt zusätzliche analytische Hypothesen.

\begin{vermutung}[EABC--Yang--Mills Continuum Bridge]
Es existiere eine eichinvariante regulierte Familie
\[
\{H_{a,L}^{\rm EABC}\}_{a,L}
\]
mit zugehörigen euklidischen Korrelationsfunktionen und quadratischen Formen
$\mathfrak q_{a,L}$.

Angenommen:

\begin{enumerate}[label=(\roman*)]
\item die regulierte Familie besitzt eine kontrollierte euklidische
Kontinuums- und unendliche-Volumen-Grenze;

\item die Grenzkorrelatoren erfüllen geeignete
Osterwalder--Schrader-Bedingungen einschließlich Reflexionspositivität;

\item der rekonstruierte physikalische Sektor definiert eine nichttriviale
vierdimensionale Yang--Mills-Theorie;

\item es existiert $\Delta_*>0$ mit der uniformen Koerzivitätsabschätzung
\[
\mathfrak q_{a,L}[u]-E_0(a,L)\|u\|^2
\ge
\Delta_*\|u\|^2
\]
für alle physikalischen
$u\perp\psi_{0,a,L}$,
uniform in $a$ und $L$ im relevanten Regime.
\end{enumerate}

Dann besitzt die rekonstruierte physikalische Theorie eine positive
Spektrallücke von mindestens $\Delta_*$.
\end{vermutung}

\begin{statusbox}{Status des Bridge Theorem Candidate}
Der letzte Implikationsschritt ist unter den formulierten Hypothesen
spektraltheoretisch unmittelbar. Die eigentliche offene Beweislast liegt in
der Konstruktion einer Yang--Mills-regulierten Familie, im kontrollierten
doppelten Grenzübergang, in der Osterwalder--Schrader-Struktur und vor allem
in der Herleitung einer \emph{uniformen} Koerzivitätskonstante $\Delta_*>0$.
Die VOA/Leech-, $H^2/H^3$- und Oktonionen-Schichten können hierfür
Strukturbedingungen liefern, ersetzen jedoch keine dieser analytischen
Hypothesen.
\end{statusbox}

\section*{7.8 Falsifikations- und Abbruchkriterien}

Die gegenwärtige EABC--Yang--Mills-Brücke scheitert in ihrer starken Form,
wenn mindestens eines der folgenden Gates verletzt wird:

\begin{enumerate}[leftmargin=1.6em]
\item Die regulierte EABC-Familie lässt sich nicht als konsistente
eichinvariante Yang--Mills-Regulierung formulieren.

\item Der Grenzübergang $a\to0$, $L\to\infty$ existiert nicht oder liefert
keine nichttriviale vierdimensionale QFT.

\item Reflexionspositivität oder eine andere notwendige
Rekonstruktionsbedingung versagt.

\item Die spektrale Lücke kollabiert:
\[
\liminf_{\substack{a\to0\\L\to\infty}}\Delta(a,L)=0.
\]

\item Die bisherige algebraische Stabilität lässt sich nicht in eine
regulatorunabhängige funktionalanalytische Abschätzung übertragen.
\end{enumerate}

Ein solcher negativer Befund widerlegt die entsprechende Yang--Mills-Brücke,
nicht jedoch automatisch die internen algebraischen EABC-Konstruktionen.

\section*{7.9 Abschließende wissenschaftliche Synthese}

Die methodische Architektur des EABC-Programms wird damit in zwei streng
getrennte Ebenen zerlegt:

\[
\boxed{
\begin{array}{c}
\textbf{Algebraisches Skelett}\\
V_4,\ Q_8,\ H^2,\ H^3,\ \text{VOA/Leech},\ \mathbb O
\end{array}}
\]

und

\[
\boxed{
\begin{array}{c}
\textbf{Analytischer Körper}\\
\text{4D-Eichfeldkonstruktion}\\
\Downarrow\\
(a,L)\text{-Grenzübergang}\\
\Downarrow\\
\text{OS-Rekonstruktion}\\
\Downarrow\\
\text{uniforme Koerzivität}\\
\Downarrow\\
\Delta_{\rm YM}>0.
\end{array}}
\]

Die erste Ebene kann notwendig sein, um einen widerspruchsfreien
diskreten Zustandsraum und geeignete Symmetrie- beziehungsweise
Sektorstrukturen zu formulieren. Sie ist jedoch nicht hinreichend für die
zweite Ebene.

\begin{warningbox}
Kein Resultat dieses Kapitels darf als Lösung oder als Teilbeweis des
Clay-Yang--Mills-Problems zitiert werden, solange weder die konstruktive
vierdimensionale Theorie noch eine regulatorunabhängige koerzitive
Spektralabschätzung für genau diese Theorie bewiesen sind.
\end{warningbox}

\end{document}
